\documentclass[14pt]{extarticle}
\usepackage{natbib}

\usepackage{xcolor}
\usepackage{amsmath}
\newcommand{\tuple}[1]{ \langle #1 \rangle }
%\usepackage{automata}
\usepackage{times}
\usepackage{ltablex}
\usepackage{tasks}
\usepackage{threeparttable}
\usepackage{booktabs}
\usepackage{xcolor} % For custom colors
\usepackage{tikz} % For styling enumerate numbers
\usepackage{tcolorbox} % For colored box styling
\usepackage{amsmath, amsfonts, amssymb, amsthm} % Math related
\usepackage{natbib}
\usepackage{fontspec}
\usepackage{luatexja}
\usepackage[mathscr]{euscript}

%%%%%% Template
\usepackage{hyperref}
\hypersetup{colorlinks=true,allcolors=blue}

\usepackage{vmargin}
\setpapersize{USletter}
\setmarginsrb{1.0in}{1.0in}{1.0in}{0.6in}{0pt}{0pt}{0pt}{0.4in}

% HOW TO USE THE ABOVE:
%\setmarginsrb{leftmargin}{topmargin}{rightmargin}{bottommargin}{headheight}{headsep}{footheight}{footskip}
%\raggedbottom
% paragraphs indent & skip:
\parindent  0.3cm
\parskip    -0.01cm

\usepackage{tikz}
\usetikzlibrary{backgrounds}

% hyphenation:
% \hyphenpenalty=10000 % no hyphen
% \exhyphenpenalty=10000 % no hyphen
\sloppy

% notes-style paragraph spacing and indentation:
\usepackage{parskip}
\setlength{\parindent}{0cm}

% let derivations break across pages
\allowdisplaybreaks

\newcommand{\orange}[1]{\textcolor{orange}{#1}}
\newcommand{\blue}[1]{\textcolor{blue}{#1}}
\newcommand{\red}[1]{\textcolor{red}{#1}}
\newcommand{\freq}[1]{{\bf \sf F}(#1)}
\newcommand{\datafreq}[2]{{{\bf \sf F}_{#1}(#2)}}

\def\qqquad{\quad\qquad}
\def\qqqquad{\qquad\qquad}

%%%%%%%%%%%%%%%%%%%%%%%%%%%%%%%%%%%%%%%%%%%%%%%%%%%%%%%%%%%%%%%%%%%%%%%%%%%%%%%%
%%%%%%%%%%%%%%%%%%%%%%%%%%%%%%%%%%%%%%%%%%%%%%%%%%%%%%%%%%%%%%%%%%%%%%%%%%%%%%%%

% fill-in-blank question style, found in https://tex.stackexchange.com/a/505089

\usepackage{ifthen}
\usepackage{tocloft}
\usepackage{exercise}
% \usepackage{xcolor}

% Set the Show Answers Boolean
\newboolean{showAns}
\setboolean{showAns}{false}
\newcommand{\showAns}{\setboolean{showAns}{true}}

% The length of the Answer line
\newlength{\answerlength}
\newcommand{\anslen}[1]{\settowidth{\answerlength}{#1}}

% ans command that indicates space for an answer or shows the answer in red
\newcommand{\ans}[1]{\settowidth{\answerlength}{\hspace{2ex}#1\hspace{2ex}}%
    \ifthenelse{\boolean{showAns}}%
        {\textcolor{red}{\underline{\hspace{2ex}#1\hspace{2ex}}}}%
        {\underline{\hspace{\answerlength}}}}%

\newcommand{\details}[1]{\settowidth{\answerlength}{#1}%
    \ifthenelse{\boolean{showAns}}%
        {\begin{tcolorbox}[colback=blue!5!white,colframe=blue!75!black,title=Solution] #1 \end{tcolorbox}}%
        {}}%


% --- Explanations for multiple choice options (shown only when \showAns is enabled) ---
\newcommand{\ExplainChoices}[4]{%
\ifthenelse{\boolean{showAns}}{%
\vspace{0.25em}
{\footnotesize
\textcolor{red}{%
\begin{itemize}
\item[(A)] #1
\item[(B)] #2
\item[(C)] #3
\item[(D)] #4
\end{itemize}}}
}{}
}

% Formatting how multiple choices Questions are formated.
\settasks{label=(\Alph*), label-width=30pt}


% Some commands for the Exercise Question package
\renewcommand{\QuestionNB}{\Large\protect\textcircled{\small\bfseries\arabic{Question}}\ }
\renewcommand{\ExerciseHeader}{} %no header
\renewcommand{\QuestionBefore}{3ex} %Space above each Q
\setlength{\QuestionIndent}{8pt} % Indent after Q number


% To create the list of answers with tocloft...
\newcommand{\listanswername}{Answers}
\newlistof[Question]{answer}{Answers}{\listanswername}

% Creates a TOC for Answers
\newcounter{prevQ}
\newcommand{\answer}[1]{\refstepcounter{answer}%
\ans{#1}%
\ifnum\theQuestion=\theprevQ%
        \addcontentsline{Answers}{answer}{\protect\numberline{}#1}% don't include the Q number
        \else%
        \addcontentsline{Answers}{answer}{\protect\numberline{\theQuestion}#1}%
        \setcounter{prevQ}{\value{Question}}%
        \fi%
        }%

% \hyphenpenalty=10000 % no hyphen
% \exhyphenpenalty=10000 % no hyphen
\sloppy              % hyphen

\newcommand{\HRule}{\rule{\linewidth}{0.5mm}}
\newcommand{\Hrule}{\rule{\linewidth}{0.3mm}}

%tocloft formatting listofanswers
\renewcommand{\cftAnswerstitlefont}{\bfseries\large}
\renewcommand{\cftanswerdotsep}{\cftnodots}
\cftpagenumbersoff{answer}
\addtolength{\cftanswernumwidth}{10pt}

\makeatletter% since there's an at-sign (@) in the command name
\renewcommand{\@maketitle}{%
  \parindent=0pt% don't indent paragraphs in the title block
  \centering
  {\Large \bfseries\textsc{\@title}} \\
  \vspace{5pt}
  {\large \textit{\@author}} \\
  \HRule \\
  \vspace{1em}
}
\makeatother% resets the meaning of the at-sign (@)


% --- Fonts (robust fallback) ---
\IfFontExistsTF{Crimson Pro Light}{
  \setmainfont{Crimson Pro Light}[
    ItalicFont={* Italic},
    BoldFont={Crimson Pro Medium},
    BoldItalicFont={Crimson Pro Medium Italic}]
  \setsansfont{Crimson Pro Light}[
    ItalicFont={* Italic},
    BoldFont={Crimson Pro Medium},
    BoldItalicFont={Crimson Pro Medium Italic}]
}{
  % HaranoAji fonts ship with TeXLive and are reliably found by LuaLaTeX
  \setmainfont{HaranoAjiMincho}
  \setsansfont{HaranoAjiGothic}
}
\title{Midterm Exam 1}
\author{Macroeconomics II \\ Hui-Jun Chen}


\begin{document}

\maketitle

% Uncomment to print answer key in red
% \showAns

\begin{Exercise}

% \section*{Instructions}
% \begin{itemize}
% \item This exam covers Lectures 1--4.
% \item \textbf{40 multiple choice questions}. Exactly one option is correct for each question.
% \item You may use a calculator. Write any scratch work on separate paper.
% \end{itemize}

% \vspace{0.8em}
% \hrule
% \vspace{0.8em}

\section*{Problem 1 (Questions 1--20): Concepts \& core mechanisms (Lectures 1--2)}

\Question \textbf{Price level.} The price level $P_t$ is best interpreted as:
\answer{B}
\begin{tasks}(1)
\task the interest rate set by the central bank
\task the value of money in terms of goods (dollars per unit of goods)
\task the real wage measured in goods
\task the quantity of money in the economy
\end{tasks}
\ExplainChoices{A policy instrument can be an interest rate, but that is not the definition of the price level.}{Correct: $P_t$ tells how many dollars are needed to buy one unit of the consumption basket (the value of money).}{The real wage is $w/P$, a different object from $P$ itself.}{$M$ is the nominal money stock; $P$ is the price of goods.}


\Question \textbf{Inflation (definition).} If we define $p_t\equiv \log P_t$, then inflation is:
\answer{C}
\begin{tasks}(1)
\task $\pi_t \equiv P_t-P_{t-1}$
\task $\pi_t \equiv \log(P_t/P_{t-1})^2$
\task $\pi_t \equiv p_t-p_{t-1}$
\task $\pi_t \equiv p_{t+1}-p_t$
\end{tasks}
\ExplainChoices{$P_t-P_{t-1}$ is a level difference, not the standard inflation definition and depends on units.}{The square is irrelevant; the usual definition is $\log(P_t/P_{t-1})$ (not squared).}{Correct: $\pi_t=p_t-p_{t-1}=\log P_t-\log P_{t-1}\approx \log(P_t/P_{t-1})$.}{$p_{t+1}-p_t$ is next period's inflation, not current inflation.}


\Question \textbf{Inflation (calculation).} If $P$ rises from $100$ to $103$ in one year, annual inflation is approximately:
\answer{B}
\begin{tasks}(1)
\task $0.3\%$
\task $3\%$
\task $30\%$
\task $-3\%$
\end{tasks}
\ExplainChoices{0.3\% would correspond to $P$ moving from 100 to 100.3, not 103.}{Correct: $(103-100)/100=0.03=3\%$ (and $\log(103/100)\approx 0.0296$).}{30\% would require 100 to 130.}{Negative inflation would require $P$ to fall.}


\Question \textbf{Nominal anchor.} A ``nominal anchor'' is:
\answer{D}
\begin{tasks}(1)
\task a policy that fixes real output in the short run
\task a rule that fixes the unemployment rate permanently
\task a mechanism that pins down the real interest rate in the long run
\task a mechanism that pins down the price level path (or long-run inflation)
\end{tasks}
\ExplainChoices{A nominal anchor does not fix real output; output is endogenously determined.}{Nominal anchor is not about permanently fixing unemployment.}{Anchoring the real rate is neither necessary nor sufficient to pin down the price level.}{Correct: a nominal anchor is what prevents $P_t$ (or long-run inflation) from drifting arbitrarily.}


\Question \textbf{Cochrane theme.} ``Deficits cause inflation'' is incomplete mainly because:
\answer{C}
\begin{tasks}(1)
\task only monetary policy can affect inflation
\task deficits never affect aggregate demand
\task inflation depends on expectations of how government debt will be backed (future surpluses)
\task inflation is purely determined by productivity
\end{tasks}
\ExplainChoices{Monetary policy affects inflation, but fiscal considerations can matter too depending on the regime.}{Deficits can boost demand; they are not irrelevant by definition.}{Correct: inflation depends on expectations of how nominal liabilities will be backed (future primary surpluses).}{Productivity affects real quantities; it does not mechanically determine the nominal price level.}


\Question \textbf{Fiscal backing intuition.} In the fiscal theory view, if expected future primary surpluses fall unexpectedly while nominal debt $B$ is given, the price level $P$ tends to:
\answer{A}
\begin{tasks}(1)
\task rise (so real debt $B/P$ falls)
\task fall (so real debt rises)
\task stay fixed because debt is nominal
\task become indeterminate only if money supply is fixed
\end{tasks}
\ExplainChoices{Correct: with nominal $B$ given, lower expected backing implies lower real value $B/P$, achieved by higher $P$.}{If $P$ fell, $B/P$ would rise, the opposite of what is needed when backing falls.}{Nominal debt being nominal is exactly why $P$ adjusts to change its real value.}{Money supply can matter in some models, but the fiscal-valuation logic does not require money-supply control.}


\Question \textbf{Money demand.} In a standard money demand function $M/P=L(Y,i)$, which sign pattern is correct?
\answer{D}
\begin{tasks}(1)
\task Increasing in income $ Y $ will decrease money demand $ L $; increasing in interest rate $ i $ will reduce money demand $ L $
\task Increasing in income $ Y $ will decrease money demand $ L $; increasing in interest rate $ i $ will increase money demand $ L $
\task Increasing in income $ Y $ will increase money demand $ L $; increasing in interest rate $ i $ will increase money demand $ L $
\task Increasing in income $ Y $ will increase money demand $ L $; increasing in interest rate $ i $ will reduce money demand $ L $
\end{tasks}
\ExplainChoices{If $L_Y<0$, money demand would fall when income rises, opposite of transactions demand.}{If $L_i>0$, people would want more money when interest rates rise, opposite of opportunity cost logic.}{If both are positive, the interest-rate effect has the wrong sign.}{Correct: higher income raises money demand ($L_Y>0$) and higher interest rates reduce it ($L_i<0$).}


\Question \textbf{LM mechanics.} Holding nominal money $M$ fixed, an increase in $P$ shifts the LM curve:
\answer{B}
\begin{tasks}(1)
\task right (lower $i$ for each $Y$)
\task left/up (higher $i$ for each $Y$)
\task does not shift because LM depends only on $M$
\task downward because inflation rises
\end{tasks}
\ExplainChoices{An LM shift right would require higher real balances, not higher $P$ with fixed $M$.}{Correct: $P\uparrow$ lowers $M/P$, shifting LM left/up (higher $i$ for each $Y$).}{LM depends on $M/P$, so changing $P$ matters.}{LM is not downward sloping in $(Y,i)$ space; it is upward sloping.}


\Question \textbf{IS slope.} In the reduced-form IS relation $Y=\bar A-a r$ with $a>0$, a higher real rate $r$ implies:
\answer{A}
\begin{tasks}(1)
\task lower $Y$ (lower demand)
\task higher $Y$ (higher saving)
\task higher inflation immediately
\task no effect on $Y$
\end{tasks}
\ExplainChoices{Correct: higher real rates reduce interest-sensitive spending, lowering equilibrium output.}{Higher saving does not raise contemporaneous output in this reduced-form demand relation.}{Inflation is not determined by IS alone in these lectures.}{IS is explicitly sensitive to $r$ through investment (and consumption in richer versions).}


\Question \textbf{Numerical IS.} Suppose $Y=\bar A-50r$ (with $r$ in decimals). If $r$ rises from $0.04$ to $0.05$, then $\Delta Y$ equals:
\answer{D}
\begin{tasks}(1)
\task $-50$
\task $-5$
\task $-0.05$
\task $-0.5$
\end{tasks}
\ExplainChoices{$-50$ would be the change if $r$ rose by 1.00 (100 percentage points).}{$-5$ corresponds to $\Delta r=0.10$, not $0.01$.}{$-0.05$ corresponds to $\Delta r=0.001$, not $0.01$.}{Correct: $\Delta r=0.01$ so $\Delta Y=-50\times 0.01=-0.5$.}


\Question \textbf{LM algebra.} Suppose money demand is $M/P=kY-hi$. Solving for $i$ gives:
\answer{A}
\begin{tasks}(1)
\task $i=\frac{k}{h}Y-\frac{1}{h}\frac{M}{P}$
\task $i=\frac{h}{k}Y-\frac{1}{k}\frac{M}{P}$
\task $i=\frac{M}{P}-kY+h$
\task $i=\frac{1}{h}\frac{M}{P}-\frac{k}{h}Y$
\end{tasks}
\ExplainChoices{Correct: rearranging gives $hi=kY-M/P$, so $i=(k/h)Y-(1/h)(M/P)$.}{This inverts the slope; units are wrong.}{Not algebraically consistent with $M/P=kY-hi$.}{Signs are reversed: higher $M/P$ should lower $i$.}


\Question \textbf{AD slope in IS--LM.} In IS--LM, why is AD downward sloping in $(P,Y)$?
\answer{C}
\begin{tasks}(1)
\task higher $P$ increases real balances $M/P$
\task higher $P$ raises potential output $Y^{FE}$
\task higher $P$ reduces $M/P$, raising $i$ and reducing demand
\task higher $P$ directly lowers productivity
\end{tasks}
\ExplainChoices{Higher $P$ reduces real balances, not increases them.}{Potential output is a supply-side concept; $P$ does not mechanically raise $Y^{FE}$.}{Correct: $P\uparrow\Rightarrow M/P\downarrow\Rightarrow i\uparrow\Rightarrow Y\downarrow$.}{Productivity is not directly affected by $P$ in IS--LM.}


\Question \textbf{Fiscal expansion.} A temporary increase in government spending $G$ (holding taxes fixed) primarily:
\answer{B}
\begin{tasks}(1)
\task shifts LM right
\task shifts IS right
\task shifts both IS and LM left
\task steepens LM without shifting it
\end{tasks}
\ExplainChoices{LM shifts with money supply or money demand, not directly with $G$.}{Correct: higher $G$ raises autonomous demand, shifting IS right.}{Both left would require contractionary policies, not higher $G$.}{Steepness depends on money demand elasticities; $G$ shifts IS.}


\Question \textbf{Monetary expansion.} An increase in nominal money $M$ primarily:
\answer{A}
\begin{tasks}(1)
\task shifts LM right/down
\task shifts IS right
\task shifts AD left
\task has no effect because $P$ adjusts instantly
\end{tasks}
\ExplainChoices{Correct: higher $M$ raises $M/P$ for given $P$, shifting LM right/down.}{$M$ does not directly shift IS in the basic setup.}{Higher $M$ tends to shift AD right, not left.}{With sticky prices in the short run, $M$ changes affect real balances and output.}


\Question \textbf{Real vs nominal (Fisher).} Which statement is correct?
\answer{B}
\begin{tasks}(1)
\task the central bank directly chooses the real rate $r$
\task the central bank chooses a nominal instrument $i$, and $r$ depends on expected inflation
\task the real rate is always equal to inflation
\task the nominal rate is independent of inflation expectations
\end{tasks}
\ExplainChoices{CB controls a nominal instrument; the real rate depends on expected inflation and inflation.}{Correct: roughly $r=i-\pi^e$; expectations matter for the real stance.}{The real rate is not equal to inflation.}{Nominal rates respond to inflation expectations through Fisher arbitrage.}


\Question \textbf{LM comparative statics.} With $M/P=kY-hi$, holding $M/P$ fixed, if $Y$ increases then $i$:
\answer{A}
\begin{tasks}(1)
\task increases
\task decreases
\task stays constant
\task becomes negative
\end{tasks}
\ExplainChoices{Correct: $\partial i/\partial Y=k/h>0$ holding $M/P$ fixed.}{If $Y$ rises, transactions demand rises, pushing rates up, not down.}{There is a mechanical relationship; it does not stay constant.}{The equation does not imply negativity generically.}


\Question \textbf{Numerical LM.} Let $M/P=6$, $k=1$, $h=2$, and $Y=8$. Then $i$ equals:
\answer{C}
\begin{tasks}(1)
\task $-1$
\task $0$
\task $1$
\task $2$
\end{tasks}
\ExplainChoices
{If you plug in the numbers, $i=(k/h)Y-(1/h)(M/P)=\tfrac{1}{2}\cdot 8-\tfrac{1}{2}\cdot 6=4-3=1$, not $-1$. So option A is wrong.}
{To get $i=0$, you would need $(k/h)Y=(1/h)(M/P)$, i.e. $Y=M/P$. With $M/P=6$, that would require $Y=6$ (not $8$).}
{Correct: $i=\tfrac{1}{2}\cdot 8-\tfrac{1}{2}\cdot 6=4-3=1$.}
{To get $i=2$, you would need $4-\tfrac{1}{2}(M/P)=2 \Rightarrow M/P=4$ (given $Y=8,k=1,h=2$), not $6$.}


\Question \textbf{Cochrane-style puzzle.} Inflation can rise even if policy rates remain near zero if:
\answer{C}
\begin{tasks}(1)
\task money demand falls only when unemployment rises
\task productivity falls sharply
\task fiscal news reduces expected backing of nominal liabilities (raising $P$)
\task central bank commits to a fixed money growth rate
\end{tasks}
\ExplainChoices{Money demand shifts are not the core here; rates can be low while inflation rises for other reasons.}{Productivity shocks affect real output; they do not directly imply inflation must rise with zero rates.}{Correct: fiscal news reducing expected backing can raise $P$ (and inflation) even before rates rise.}{Fixed money growth is one regime, but the question targets the near-zero-rate episode logic.}


\Question \textbf{Deficits and inflation (evidence logic).} ``Large deficits are not always followed by inflation'' is consistent with:
\answer{A}
\begin{tasks}(1)
\task deficits expected to be offset by future surpluses (credible backing)
\task deficits always being fully monetized
\task deficits always lowering real rates permanently
\task inflation being unrelated to expectations
\end{tasks}
\ExplainChoices{Correct: if deficits are expected to be reversed by future surpluses, inflation need not rise much.}{Full monetization is precisely the case that tends to create inflation pressure.}{Real rates can rise or fall depending on policy; deficits do not guarantee permanent rate declines.}{Expectations are central in modern inflation analysis.}


\Question \textbf{Anchoring (concept).} Which situation best reflects an \emph{unanchored} nominal anchor?
\answer{D}
\begin{tasks}(1)
\task long-run expected inflation stays near the target after shocks
\task inflation returns quickly to target after a temporary supply shock
\task the central bank reacts more when inflation is above target
\task private agents revise long-run expected inflation upward after a temporary shock
\end{tasks}
\ExplainChoices{Anchored expectations means this stays near target, not unanchored.}{Quick return to target suggests anchoring, not unanchoring.}{A stronger reaction can help anchoring; it is not itself unanchoring.}{Correct: unanchored means agents revise long-run expected inflation upward after shocks.}


\vspace{0.8em}
\hrule
\vspace{0.8em}

\section*{Problem 2 (Questions 21--40): AD--AS and IS--MP math (Lectures 3--4)}

\Question \textbf{Supply shock in AD--AS.} A negative supply shock typically causes:
\answer{B}
\begin{tasks}(1)
\task $Y\uparrow$ and $P\downarrow$
\task $Y\downarrow$ and $P\uparrow$
\task $Y\uparrow$ and $P\uparrow$
\task $Y\downarrow$ and $P\downarrow$
\end{tasks}
\ExplainChoices{This is not the typical supply-shock pattern; supply shocks reduce output.}{Correct: adverse supply raises costs, shifting AS up/left: $Y\downarrow$, $P\uparrow$.}{$Y\uparrow$ typically comes from demand shifts, not negative supply.}{Both down is typical of negative demand shocks.}


\Question \textbf{Demand shock in AD--AS.} A positive demand shock typically causes:
\answer{C}
\begin{tasks}(1)
\task $Y\downarrow$ and $P\uparrow$
\task $Y\downarrow$ and $P\downarrow$
\task $Y\uparrow$ and $P\uparrow$
\task $Y\uparrow$ and $P\downarrow$
\end{tasks}
\ExplainChoices{$Y\downarrow$, $P\uparrow$ is the supply-shock pattern.}{Both down is negative demand.}{Correct: positive demand shifts AD right: $Y\uparrow$, $P\uparrow$.}{$P\downarrow$ would require AS improvements or negative demand.}


\Question \textbf{Full-employment output.} In your Lecture 3 notation, $Y^{FE}$ is:
\answer{A}
\begin{tasks}(1)
\task potential output determined by the RBC supply block
\task output chosen by the central bank
\task output when inflation is zero
\task output when money supply is constant
\end{tasks}
\ExplainChoices{Correct: $Y^{FE}$ is potential output pinned by technology and market clearing in the RBC block.}{CB sets policy instruments, not potential output.}{Inflation can be zero at many output levels depending on expectations/shocks.}{$Y^{FE}$ is not tied to money growth mechanically.}


\Question \textbf{Medium run.} In the medium run of the AD--AS story, output tends to return toward $Y^{FE}$ mainly because:
\answer{D}
\begin{tasks}(1)
\task AD shifts back automatically
\task productivity always rises
\task money supply returns to its initial level
\task expected prices (or wages) adjust, shifting SRAS
\end{tasks}
\ExplainChoices{AD need not shift back automatically; it depends on policy and shocks.}{Productivity is not assumed to always rise.}{Money supply may not return; modern policy targets rates.}{Correct: as wages/prices/expectations adjust, SRAS shifts and output returns toward potential.}


\Question \textbf{SRAS (level form).} If $y=y^{FE}+\alpha(p-p^e)$, then the output gap is:
\answer{B}
\begin{tasks}(1)
\task $x=y+y^{FE}$
\task $x=\alpha(p-p^e)$
\task $x=\alpha(p^e-p)$
\task $x=p-p^e$
\end{tasks}
\ExplainChoices{$x=y-y^{FE}$, not $y+y^{FE}$.}{Correct: $x=y-y^{FE}=\alpha(p-p^e)$ from the given SRAS.}{This flips the sign of the surprise term.}{$p-p^e$ misses the slope parameter $\alpha$ and ignores $y^{FE}$.}


\Question \textbf{Okun's law sign.} If $u-u^n\approx -\gamma x$ with $\gamma>0$, then $x>0$ implies:
\answer{A}
\begin{tasks}(1)
\task unemployment below natural ($u<u^n$)
\task unemployment above natural ($u>u^n$)
\task unemployment unchanged
\task inflation must be falling
\end{tasks}
\ExplainChoices{Correct: $x>0$ means output above potential, so unemployment is below natural.}{$u>u^n$ corresponds to slack ($x<0$).}{It changes unless $x=0$.}{Inflation dynamics requires a Phillips curve; Okun alone does not pin it down.}


\Question \textbf{Numerical Okun.} If $\gamma=0.5$ and the output gap is $x=2\%$, the unemployment gap is approximately:
\answer{C}
\begin{tasks}(1)
\task $+4\%$
\task $+1\%$
\task $-1\%$
\task $-4\%$
\end{tasks}
\ExplainChoices{+4\% has the wrong sign and magnitude.}{+1\% has the wrong sign.}{Correct: $u-u^n\approx -0.5\times 2\%=-1\%$.}{$-4\%$ is too large in magnitude.}


\Question \textbf{Log utility Euler (calculation).} With $C_{t+1}/C_t=\beta(1+r_t)$, if $\beta=0.96$ and $r_t=0.05$, then $C_{t+1}/C_t$ equals:
\answer{B}
\begin{tasks}(1)
\task $0.048$
\task $1.008$
\task $0.992$
\task $1.09375$
\end{tasks}
\ExplainChoices{0.048 is $\beta r$, not $\beta(1+r)$.}{Correct: $0.96\times 1.05=1.008$.}{0.992 is the reciprocal, not the product.}{1.09375 is $1.05/0.96$, not $0.96\cdot 1.05$.}


\Question \textbf{Steady-state real rate.} Under $1=\beta(1+\bar r)$ with $\beta=0.96$, $\bar r$ is approximately:
\answer{D}
\begin{tasks}(1)
\task $0.96$
\task $0.04$
\task $-0.04$
\task $1/0.96-1\approx 0.0417$
\end{tasks}
\ExplainChoices{$\beta$ itself is not a rate.}{0.04 is close but missing the exact calculation.}{$-0.04$ has the wrong sign.}{Correct: $\bar r=1/\beta-1\approx 0.0417$.}


\Question \textbf{Dynamic IS (log utility).} A standard IS in gap form is:
\answer{A}
\begin{tasks}(1)
\task $x_t=\mathbb{E}_t x_{t+1}-(r_t-r_t^n)+\varepsilon_t^d$
\task $x_t=\mathbb{E}_t x_{t+1}+(r_t-r_t^n)+\varepsilon_t^d$
\task $x_t=(r_t-r_t^n)+\varepsilon_t^d$
\task $x_t=-(r_t-r_t^n)-\varepsilon_t^d$
\end{tasks}
\ExplainChoices{Correct: with log utility, coefficient on $(r-r^n)$ is 1 and it enters negatively.}{Wrong sign: higher $r$ would raise demand here.}{Omits expectations term and has wrong structure.}{Flips both signs.}


\Question \textbf{Real-rate MP rule.} A real-rate policy rule consistent with Lecture 4 is:
\answer{B}
\begin{tasks}(1)
\task $r=\bar r-\phi_\pi(\pi-\pi^*)-\phi_y(Y-Y^{FE})$
\task $r=\bar r+\phi_\pi(\pi-\pi^*)+\phi_y(Y-Y^{FE})$
\task $r=\pi+\phi_\pi(\pi-\pi^*)$
\task $r=\bar r$ always
\end{tasks}
\ExplainChoices{The signs correspond to a destabilizing response to inflation/overheating.}{Correct: higher inflation or overheating triggers a higher real-rate stance.}{Not a standard policy rule; omits intercept/gap response appropriately.}{A constant real rate is not inflation targeting.}


\Question \textbf{IS--MP implication.} In $Y=\bar A-a r$ and $r=\bar r+\phi_\pi(\pi-\pi^*)$, higher inflation implies:
\answer{D}
\begin{tasks}(1)
\task higher $Y$ because $r$ falls
\task no change in $Y$ because $r$ is real
\task higher $Y$ because policy accommodates inflation
\task lower $Y$ because policy raises $r$
\end{tasks}
\ExplainChoices{This would imply policy cuts $r$ when inflation rises, which would fuel inflation.}{Real policy can change output via IS; it is not irrelevant.}{Accommodation is the opposite of the stated rule.}{Correct: inflation above target raises $r$, lowering $Y$ through IS.}


\Question \textbf{Numerical policy.} Suppose $r=0.02+2(\pi-\pi^*)$ with $\pi^*=0.02$. If inflation rises from $2\%$ to $3\%$, then $r$ rises by:
\answer{A}
\begin{tasks}(1)
\task $0.02$
\task $0.01$
\task $0.002$
\task $0.2$
\end{tasks}
\ExplainChoices{Correct: $\Delta\pi=0.01$, coefficient 2 gives $\Delta r=0.02$.}{0.01 would be coefficient 1.}{0.002 is too small by a factor of 10.}{0.2 is too large by a factor of 10.}


\Question \textbf{Numerical IS--MP.} If $Y=\bar A-10r$ and $r$ rises by $0.02$, then $\Delta Y$ equals:
\answer{C}
\begin{tasks}(1)
\task $+0.2$
\task $0$
\task $-0.2$
\task $-2$
\end{tasks}
\ExplainChoices{+0.2 would require $r$ to fall, not rise.}{Output changes because $r$ changes in IS.}{Correct: $\Delta Y=-10\times 0.02=-0.2$.}{$-2$ would require $\Delta r=0.2$.}


\Question \textbf{AD--AS solving (comparative statics).} With AD $y=\bar y_D-\phi p$ and SRAS $p=p^e+\lambda(y-y^{FE})$ (holding $p^e$ fixed), a higher $\bar y_D$ leads to:
\answer{B}
\begin{tasks}(1)
\task lower $p$ and lower $y$
\task higher $p$ and higher $y$
\task higher $p$ and lower $y$
\task no change in equilibrium
\end{tasks}
\ExplainChoices{This would correspond to a negative demand shift.}{Correct: AD shifts right, so both $y$ and $p$ rise when SRAS slopes upward.}{Price up with output down is supply contraction.}{Equilibrium changes unless slopes are degenerate.}


\Question \textbf{Solve once (formula).} From AD $y=\bar y_D-\phi p$ and SRAS $p=p^e+\lambda(y-y^{FE})$, the equilibrium price level satisfies:
\answer{D}
\begin{tasks}(1)
\task $p=p^e+\lambda(\bar y_D-y^{FE})$
\task $p=\dfrac{p^e}{1+\lambda\phi}$
\task $p=\dfrac{\bar y_D-y^{FE}}{1+\lambda\phi}$
\task $p=\dfrac{p^e+\lambda(\bar y_D-y^{FE})}{1+\lambda\phi}$
\end{tasks}
\ExplainChoices{Ignores the AD slope interaction; you must solve simultaneously.}{Missing the demand shifter term.}{Ignores $p^e$ and $\lambda$.}{Correct: $(1+\lambda\phi)p=p^e+\lambda(\bar y_D-y^{FE})$.}


\Question \textbf{Policy aggressiveness.} In an IS--MP AD derivation, increasing $\phi_\pi$ (holding other parameters fixed) generally makes AD:
\answer{A}
\begin{tasks}(1)
\task steeper (output responds more to inflation)
\task flatter (output responds less to inflation)
\task shift right
\task shift left
\end{tasks}
\ExplainChoices{Correct: stronger inflation response means a larger output contraction for a given inflation increase (steeper AD).}{Flatter would correspond to weaker response.}{Steepness changes, not a parallel shift.}{Steepness changes, not a parallel shift.}


\Question \textbf{Monetary easing (IS--MP).} In the IS--MP view, an ``AD shift right due to easier monetary policy'' is best interpreted as:
\answer{D}
\begin{tasks}(1)
\task an increase in $M$ holding the policy rule fixed
\task a fall in productivity that raises $Y^{FE}$
\task an upward shift in SRAS due to cost shocks
\task a lower intercept $\bar r$ / more accommodative stance in the real-rate rule
\end{tasks}
\ExplainChoices{In IS--MP we do not interpret easing primarily as $M\uparrow$.}{This is a supply-side change, not monetary easing.}{This is cost-push, not easing.}{Correct: easing corresponds to a lower intercept/more accommodative stance (lower $\bar r$).}


\Question \textbf{Cochrane connection.} In a regime with uncertain fiscal backing, an interest-rate rule alone may not pin down the price level because:
\answer{B}
\begin{tasks}(1)
\task the IS curve becomes vertical
\task nominal liabilities must be valued by expected surpluses; without backing, $P$ adjusts through valuation
\task unemployment becomes constant
\task money demand becomes infinite
\end{tasks}
\ExplainChoices{A vertical IS is not the key issue in Cochrane's regime discussion.}{Correct: without credible fiscal backing, $P$ can adjust via valuation of nominal liabilities; rate rules alone may not anchor $P$.}{Unemployment is not fixed by this argument.}{Money demand need not be infinite.}


\Question \textbf{Level vs rate.} A permanent one-time increase in the price level $P$ with no further changes implies:
\answer{C}
\begin{tasks}(1)
\task permanently high inflation
\task permanently negative inflation
\task a one-time burst of inflation, then inflation returns to normal
\task inflation is undefined
\end{tasks}
\ExplainChoices{Permanent high inflation means $\pi$ stays high period after period.}{Permanent deflation means $P$ keeps falling.}{Correct: a one-time jump in $P$ implies a one-time inflation spike, then inflation can return to normal.}{Inflation is well-defined as the growth rate of $P$.}


\end{Exercise}

\end{document}
